\sourcepage{368}
\chapterappendix{11A}{预试验}
\sourcepage{369}
许多研究，尤其大型队列和临床试验，可以从主研究开始前、甚至大研究资助申请撰写前开展的预试验中，获得有益经验和资料。预试验常是了解干预类型、剂量与持续时间，招募、随机化及保留参与者的可行性，特定测量的障碍，依从可能性，潜在不良事件及研究费用的最佳方法。其范围可以从少量志愿者中短暂试行一种测量，到数百人参加、为多中心多年研究作准备的较长期研究。

有些预试验的重要目标，是在尽量减少不良事件机会的同时，确定临床试验的最佳干预。例如，可比较数种剂量对中间标志物的作用。如果属于申请FDA批准的系列研究，通常称为II期试验。如果目的仅限于证明计划中的测量、资料收集工具及数据管理系统可行，估计主试验成本，以及取得连续结局变化的资料，就可能不需要对照组。

预试验有时用于估计样本量参数，例如安慰剂组的结局发生率或均数，预期干预效应，以及统计变异程度。但多数情况下，最好从可比参与者、类似干预的已发表研究取得这些估计；没有这些资料时，预试验才可能有帮助。尤其是将作为主要终点的连续测量变化的标准差，既往研究往往没有提供，可在一组类似主试验预期参与者的小样本中，随时间重复测量取得。不过，预试验不是预期效应量的可靠来源：\sourcepage{370}如果它已有足够效能精确确定效应，就已经实现主试验的目的了！

好的预试验需要相当的时间和资源，但也提高主试验获资助并顺利完成的机会。大型试验的预试验应有简短而完整、经IRB批准的方案、资料收集表和分析计划；这些准备可节约主试验启动时间。除正式试验预期的预测因素和结局外，还可测量可供招募的人数、对不同招募方法响应的比例、合格但可能拒绝随机的比例、招募及随机化所需时间和费用，以及对干预、访视和方案其他方面的依从程度。结束后听取参与者与工作人员对试验方法如何改进的意见，通常很有帮助。

\chapterappendix{11B}{三项提前停止的试验}
\sourcepage{371}
下面介绍导致三项试验提前停止的资料，以及提前终止产生的影响。

\subsection{耐甲氧西林金黄色葡萄球菌菌血症联合抗生素：因害停止（1）}
这项开放标签随机试验比较住院耐甲氧西林金黄色葡萄球菌（methicillin-resistant \textit{Staphylococcus aureus}，MRSA）菌血症的两种抗生素方案，判断标准治疗加入抗葡萄球菌β内酰胺，是否优于单用标准治疗。研究者计划随机440人，检出死亡、5天后持续菌血症、微生物学复发及微生物学治疗失败复合结局的绝对风险降低12.5\%。入组并随访343人后，中期分析显示主要结局为联合组35\%、标准组39\%，差$-4\%$（95\%置信区间$-14\%$～6\%，$P=0.42$），没有显著差异。七项预定次要结局中，包括全因死亡在内的五项也无显著差异。但急性肾损伤在联合组更常见，分别23\%和6\%，差17\%（95\%置信区间9\%～25\%）；第5天持续菌血症则较少，分别11\%和20\%，差$-9\%$（$-17\%$～$-1\%$）。

DSMB在此前220人的中期分析中，已首次注意到肾损伤差异；增加123人后，联合组肾损伤仍较多，90天死亡却未减少，因此建议停止招募。\sourcepage{372}研究者解释，持续菌血症减少虽可视为联合方案更有效的依据，但他们认为其收益被肾损伤的潜在危害抵消。不过，样本量也可能不足以发现其他有利于联合方案、具有临床重要性的差异。

\subsection{早期乳腺癌他莫昔芬后用来曲唑：因益停止（2）}
研究开展时，已知早期激素受体阳性乳腺癌女性服雌激素调节剂他莫昔芬5年，可提高无病生存和总体生存。研究者开展双盲、随机、安慰剂对照试验，判断完成他莫昔芬后，再用非甾体芳香化酶抑制剂来曲唑（letrozole）5年，是否进一步提高无病生存。首次中期分析时已入5157人，中位随访仅2.4年；来曲唑组75人、安慰剂组132人发生复发或对侧乳房新原发癌，4年无病生存率分别93\%和87\%（$P<0.001$）。DSMB据此建议公开结果、告知治疗分配，并给安慰剂者接受来曲唑的机会，符合原统计监测及早停规则。但有人认为，中位仅2.4年、尚无人随访满5年就停止，降低了结果的有用性（3）。它也限制了长期不良作用的研究。停止时，来曲唑组自报骨质疏松诊断较多的趋势已出现（$P=0.07$），却无法再取得更明确的资料。

\subsection{胰腺切除常规腹膜内引流：因害停止一个亚组（4）}
\sourcepage{373}
这项试验研究胰腺切除后常规放置腹膜内引流管的做法。许多外科医师用它排出血液或胰液，但也有人担心增加严重并发症，例如胰瘘，即液体经不愈合的引流伤口通道持续排出。研究者开展未盲随机试验，考察引流对成人术后并发症频率和严重程度的影响，包括胰十二指肠切除（pancreaticoduodenectomy，PD，切除胰头）和远端胰切除（切除胰体或胰尾）。计划从美国9个学术外科中心招募750人，检出两组2级或以上并发症发生率10\%的差异；此类并发症可能危及生命，并需要干预或延长住院。

但招募达到目标18\%时的中期监测（$N=282$，含196例PD），发现PD无引流者死亡8人（12\%），有引流者2人（3\%），$P=0.10$。\footnote{译注：底本同时写计划750人、已入282人和“目标18\%”；282／750并非18\%。此处保留全部原数，不代改百分比或样本数。}DSMB建议停止PD亚组，继续招募远端胰切除者。其他临床科学家因死亡数很少而对决定感到意外，并要求说明中期监测计划。研究者解释，原先确有针对主要结局的统计早停计划，却没有监测死亡的计划（5）。但DSMB已知一项新发表、超过千人胰切除的回顾分析，报告无腹膜内引流的PD患者死亡较多。由于本试验也有死亡增加趋势，研究者加入了针对死亡的停止计划。\sourcepage{374}最终完成随访的344例远端胰切除者，死亡或2级以上并发症发生率均未见差异（6）。

\begin{references}
\item Ton SY, Lye DC, Yahav D, et al. Effect of vancomycin or daptomycin with vs without an antistaphylococcal β-lactam on mortality, bacteremia, relapse, or treatment failure in patients with MRSA bacteremia: a randomized clinical trial. JAMA. 2020;323(6):527-537.
\item Goss PE, Ingle JN, Martino S, et al. A randomized trial of letrozole in postmenopausal women after five years of tamoxifen therapy for early-stage breast cancer. N Engl J Med. 2003;349(19):1793-1802.
\item Bryan J, Wormack N. Letrozole after tamoxifen for breast cancer—what is the price of success? N Engl J Med. 2003;349(19):1855-1857.
\item Van Buren G 2nd, Bloomston M, Hughes SJ, et al. A randomized prospective multicenter trial of pancreaticoduodenectomy with and without routine intraperitoneal drainage. Ann Surg. 2014;259(4):605-612.
\item Van Bruen G 2nd, Fisher WE. Pancreaticoduodenectomy without drains: interpretation of the evidence. Ann Surg. 2016;263(2):e20-e21.
\item Van Bruen G 2nd, Bloomston M Schmidt CR, et al. A prospective randomized multicenter trial of distal pancreatectomy with and without routine intraperitoneal drainage. Ann Surg. 2017;266(3):421-431.
\end{references}


\chapterappendix{11C}{第11章习题：随机盲法试验的设计}
\sourcepage{375}
\begin{enumerate}
\item 草药提取物石杉碱（huperzine）在中国曾用于治疗痴呆，动物及人体初步研究结果令人期待。你希望检验它能否减缓阿尔茨海默病进展。研究发现，血浆淀粉样β蛋白1—40（Aβ40）是阿尔茨海默病的生物标志物：其升高与痴呆发生风险显著增加有关，并随痴呆进展而增加。你规划在轻度认知障碍老年患者中检验石杉碱预防痴呆的疗效，考虑两种结局：Aβ40水平变化，或新出现临床痴呆诊断。
\begin{enumerate}
\item 将Aβ40水平变化作为主要结局，列出一项优点和一项缺点。
\item 将新诊断痴呆作为主要结局，列出一项优点和一项缺点。
\end{enumerate}
\item 你的石杉碱试验主要检验：它是否降低有轻度认知障碍的老年男女出现临床痴呆诊断的发生率。为实现下列各目标，说明应在基线收集哪类信息，并分别列出两个潜在变量：
\begin{enumerate}
\item 尽可能提高团队随访和保留参与者的能力；
\item 使临床医师和其他研究者评估样本的可推广性；
\item 确定石杉碱能否有效减缓阿尔茨海默病进展；
\item 日后通过亚组分析，评估某些人群中的潜在治疗作用差异；
\item 使资料还能回答推动本试验的主要问题之外的其他问题。
\end{enumerate}
\item 携带ApoE4等位基因者痴呆风险较高。你怀疑石杉碱在携带与未携带者中的作用不同。
\begin{enumerate}
\item 你会如何设计（1）筛查和入组程序，以及（2）随机化策略，以回答该问题？\sourcepage{376}
\item 在石杉碱的初次试验中尝试回答这个问题，有哪些潜在好处和弊端？
\end{enumerate}
\item 目前正计划一项较大的石杉碱随机安慰剂对照试验，随访2年。药物可能引起腹泻、恶心和呕吐等胃肠症状。下列哪一种不良作用监测方法更好，为什么？
\begin{enumerate}
\item 结束时询问，研究治疗是否使参与者出现腹泻、恶心或呕吐。
\item 开始时告知参与者可随时报告任何可能的副作用，然后等待他们主动报告胃肠症状。
\item 每次随访都用包括腹泻、恶心、呕吐的清单，询问是否出现这些症状。
\item 每次随访问开放式问题，例如：“上次访视以来，你有任何新症状或健康状况吗？”
\end{enumerate}
\item 部分参与者因胃肠症状停用石杉碱。作为中期监察者，你正考虑是否因为胃不适、恶心或呕吐报告较多，建议提前结束试验。除了发生频率，这些有症状者的哪些信息有助于决定？除提前终止外，还可考虑哪些做法？
\item 试验期间，20\%随机参与者未回来参加1年访视，40\%在2年内停止研究药物。严格按意向治疗分析治疗对胃肠副作用的影响，列出一项优点和一项缺点。
\item 2年意向治疗分析显示，石杉碱组比安慰剂组较少被诊断为痴呆，但差异不显著（$P=0.08$）。另一分析提示，在60岁以下者，石杉碱比安慰剂多降低25\%的痴呆诊断可能性，该亚组$P=0.01$。据此断言石杉碱可有效预防60岁以下者临床痴呆，可能有哪些问题？
\end{enumerate}
